\chapter{Skojarzenia, systemy reprezentantów i cykle Hamiltona}
\label{ch:tg-skojarzenia}

Drugim podstawowym problemem kombinatorycznym jest wyszukiwanie
wielu krawędzi, które nie mogą ze sobą kolidować. Skojarzenia
opisują przydziały, równoległe parowania i układy niezależnych
wyborów. Na końcu rozdziału wrócimy do przechodzenia po grafach,
ale zamiast wszystkich \emph{krawędzi} spróbujemy odwiedzić
wszystkie \emph{wierzchołki} dokładnie raz. Takie jest
zasadnicze rozróżnienie między teorią Eulera i Hamiltona.

\section{Skojarzenia i ścieżki powiększające}
\label{sec:tg-matchings}

\begin{definicja}[Skojarzenie]
Niech $G=(V,E)$ będzie skończonym grafem prostym.
\emph{Skojarzeniem} nazywamy podzbiór $M\subseteq E$,
w którym żadne dwie krawędzie nie mają wspólnego końca.
Wierzchołek jest \emph{skojarzony}, jeśli jest końcem
pewnej krawędzi z $M$, a w przeciwnym razie jest
\emph{wolny}. Skojarzenie jest \emph{maksymalne}, gdy
nie można dodać do niego kolejnej krawędzi bez
naruszenia warunku; jest \emph{największe}, gdy ma
największą możliwą liczbę krawędzi. Jest \emph{doskonałe},
gdy każdy wierzchołek jest skojarzony. Dopuszczamy
skojarzenie puste; w grafie zerowym jest ono doskonałe,
a w grafie o nieparzystej liczbie wierzchołków
skojarzenie doskonałe nie istnieje, ponieważ każda
krawędź skojarzenia pokrywa dokładnie dwa różne wierzchołki.
\end{definicja}

Maksymalność nie oznacza największej liczności.
Przykładowo w ścieżce czterowierzchołkowej $a-b-c-d$
skojarzenie $\{bc\}$ jest maksymalne, lecz nie największe:
$\{ab,cd\}$ ma dwie krawędzie.

\begin{definicja}[Ścieżka naprzemienna i powiększająca]
Niech $G=(V,E)$ będzie skończonym grafem prostym,
a $M$ jego skojarzeniem. \emph{Ścieżką
naprzemienną względem $M$} nazywamy ścieżkę, na której
krawędzie należące do $M$ i do $E\setminus M$
występują naprzemiennie. Jest ona \emph{powiększająca},
jeżeli oba jej końce są wolne, a pierwsza i ostatnia
krawędź nie należą do $M$. Dopuszczamy tylko ścieżki
z co najmniej jedną krawędzią.
\end{definicja}

\begin{lemat}[Odwrócenie ścieżki powiększającej]
Niech $P$ będzie ścieżką powiększającą względem
skojarzenia $M$ w grafie $G$. Wtedy
$M':=M\mathbin{\triangle}E(P)$ jest skojarzeniem,
a $|M'|=|M|+1$, gdzie $\triangle$ oznacza różnicę
symetryczną zbiorów.
\end{lemat}

\begin{proof}
Na wewnętrznych wierzchołkach $P$ jedna z dwóch
krawędzi ścieżki należy do $M$, a druga nie.
Po zamianie ról tych krawędzi każdy taki wierzchołek
nadal jest incydentny dokładnie z jedną krawędzią
nowego skojarzenia. Żadna krawędź $M$ spoza $E(P)$
nie może być incydentna z wewnętrznym wierzchołkiem
$P$, ponieważ byłaby drugą krawędzią skojarzenia
przy tym wierzchołku. Końce $P$ były wolne,
a teraz dostają po jednej krawędzi. Poza $P$
skojarzenie się nie zmienia. Ścieżka rozpoczyna się
i kończy krawędzią spoza $M$, zatem liczba usuniętych
krawędzi $M$ jest o jeden mniejsza od liczby dodanych.
Daje to $|M'|=|M|+1$.
\end{proof}

\begin{center}
\begin{tikzpicture}[scale=.85,line cap=round]
\foreach \i/\n in {0/a,1/b,2/c,3/d,4/e,5/f}
 {\node[tgVertex] (\n) at (1.3*\i,0) {$\n$};}
\draw[tgTree,line width=2.6pt] (b)--(c) (d)--(e);
\draw[tgAccent,line width=2.5pt] (a)--(b) (c)--(d) (e)--(f);
\node[font=\footnotesize,text=tgNavy] at (3.25,-.80)
 {Zielone krawędzie należą do $M$, czerwone — do $E\setminus M$.};
\end{tikzpicture}

{\footnotesize Po zamianie kolorów na tej ścieżce
otrzymujemy skojarzenie o jedną krawędź większe.}
\end{center}

\begin{twierdzenie}[Berge]
\label{thm:tg-berge}
Skojarzenie $M$ skończonego grafu prostego $G$
jest największe wtedy i tylko wtedy, gdy
nie istnieje ścieżka powiększająca względem $M$.
\end{twierdzenie}

\begin{proof}
Jeśli istnieje ścieżka powiększająca, poprzedni lemat
daje skojarzenie większe od $M$.

Załóżmy odwrotnie, że $M$ nie jest największe.
Wybieramy skojarzenie $N$ o $|N|>|M|$.
Rozważmy graf $H=(V(G),M\triangle N)$.
Ponieważ każdemu wierzchołkowi przylega co najwyżej
jedna krawędź z $M$ i co najwyżej jedna z $N$,
w grafie tym stopień każdego wierzchołka jest co
najwyżej dwa. Uzasadnijmy, że każda nieizolowana
składowa $H$ jest ścieżką albo cyklem. Jeśli zawiera
wierzchołek stopnia jeden, zaczynamy w nim i idziemy
kolejnymi krawędziami, nie wracając poprzednią krawędzią.
Nie możemy ponownie odwiedzić początku, gdyż ma on
stopień jeden, ani innego wcześniej odwiedzonego
wierzchołka, gdyż wymagałoby to trzeciej krawędzi
incydentnej z nim. Ze skończoności grafu wynika,
że kończymy w innym wierzchołku stopnia jeden.
Wewnętrzne wierzchołki ścieżki mają już po dwie
krawędzie, a końce po jednej i nie mają innych sąsiadów.
Spójność składowej oznacza więc, że jest ona tą ścieżką.
Jeśli składowa nie zawiera wierzchołka stopnia jeden,
wszystkie jej wierzchołki mają stopień dwa.
Idąc analogicznie od dowolnego wierzchołka, musimy
powrócić do początku (powrót do innego wierzchołka
wymagałby trzeciej krawędzi). Uzyskany cykl wykorzystuje
oba końce incydencji w każdym swoim wierzchołku,
więc ze spójności jest całą składową. Ponieważ na każdej składowej
kolejne krawędzie muszą na przemian należeć do $M$
i $N$, każdy cykl ma parzystą długość i jednakową
liczbę krawędzi z obu skojarzeń, a więc nie zmienia różnicy ich
liczności. Ścieżki mają równe liczby krawędzi
obu typów albo różnią się o jedną.

Ponieważ $|N\setminus M|-|M\setminus N|=|N|-|M|>0$,
suma różnic liczebności obu typów krawędzi po
składowych $H$ jest dodatnia. Istnieje więc składowa
będąca ścieżką z jedną krawędzią $N$ więcej
niż $M$. Jej pierwsza i ostatnia krawędź
należą do $N\setminus M$. Jej końce są wolne
w $M$: gdyby któryś był skojarzony przez $M$,
odpowiednia krawędź $e\in M$ nie mogłaby należeć do $N$:
w $N$ koniec ścieżki jest już pokryty inną krawędzią
z $N\setminus M$. Wtedy $e\in M\setminus N$
i koniec ścieżki miałby w $H$ stopień co najmniej dwa,
wbrew temu, że jest końcem składowej będącej ścieżką. Otrzymaliśmy ścieżkę
powiększającą względem $M$.
\end{proof}

\section{Grafy dwudzielne i twierdzenie Halla}
\label{sec:tg-hall}

\begin{definicja}[Sąsiedztwo zbioru]
Niech $G=(A\sqcup B,E)$ będzie skończonym grafem
dwudzielnym o ustalonych częściach $A$ i $B$.
Dla $S\subseteq A$ kładziemy
$N(S):=\{b\in B:\exists a\in S,\ ab\in E\}$.
Mówimy, że skojarzenie \emph{pokrywa $A$}, jeżeli
każdy wierzchołek $a\in A$ jest końcem pewnej
krawędzi skojarzenia.
\end{definicja}

\begin{twierdzenie}[Hall: twierdzenie o małżeństwach]
\label{thm:tg-hall}
Niech $G=(A\sqcup B,E)$ będzie skończonym grafem
dwudzielnym. Istnieje skojarzenie pokrywające
wszystkie wierzchołki zbioru $A$ wtedy i tylko wtedy,
gdy dla każdego $S\subseteq A$ zachodzi
\[
|N(S)|\ge |S|.
\]
\end{twierdzenie}

\begin{proof}[Dowód indukcyjny]
Konieczność jest natychmiastowa: różne wierzchołki
$S$ muszą być skojarzone z różnymi wierzchołkami
zbioru $N(S)$.

Dla dostateczności indukujemy względem $|A|$.
Jeśli $|A|=0$, skojarzenie puste działa.
Jeśli $|A|=1$, warunek Halla daje choć jednego
sąsiada, więc możemy wybrać jedną krawędź.
Załóżmy, że $|A|\ge2$ i twierdzenie zostało
udowodnione dla mniejszych części lewych.

\emph{Przypadek pierwszy:} istnieje niepusty, właściwy
zbiór $S\subsetneq A$ z $|N(S)|=|S|$.
W podgrafie indukowanym na $S\sqcup N(S)$
sąsiedztwo każdego $R\subseteq S$ jest nadal $N(R)$,
bo $N(R)\subseteq N(S)$. Warunek Halla jest więc
spełniony i z założenia
indukcyjnego istnieje skojarzenie pokrywające $S$.
Na grafie z częściami $A\setminus S$
i $B\setminus N(S)$ warunek również jest spełniony:
dla $T\subseteq A\setminus S$ mamy
$|N(S\cup T)|\ge|S|+|T|$. Ponieważ
$N(S\cup T)=N(S)\cup N(T)$ oraz $|N(S)|=|S|$,
otrzymujemy $|N(T)\setminus N(S)|\ge|T|$.
Lewa strona jest mocą sąsiedztwa $T$ w pozostałym
grafie indukowanym.
Stąd drugie skojarzenie pokrywa $A\setminus S$.
Ponieważ zbiory wierzchołków użytych przez obydwa
skojarzenia są rozłączne, ich suma pokrywa $A$.

\emph{Przypadek drugi:} dla każdego niepustego
właściwego $S\subsetneq A$ zachodzi
$|N(S)|\ge|S|+1$. Wybierzmy dowolne $a\in A$.
Warunek Halla dla $\{a\}$ zapewnia istnienie
sąsiada $b\in B$. Usuwamy $a$ i $b$.
Dla dowolnego niepustego $T\subseteq A\setminus\{a\}$
zbiór $T$ jest właściwym podzbiorem $A$, więc
$|N(T)\setminus\{b\}|\ge|N(T)|-1\ge|T|$.
Dla $T=\varnothing$ warunek jest automatyczny.
Z założenia indukcyjnego otrzymujemy skojarzenie
pokrywające $A\setminus\{a\}$ w pozostałym grafie;
po dodaniu krawędzi $ab$ uzyskujemy skojarzenie
pokrywające $A$.
Oba przypadki wyczerpują możliwości.
\end{proof}

\begin{przyklad}[Przydział trzech osób do trzech stanowisk]
Niech $A=\{a_1,a_2,a_3\}$ oznacza osoby,
$B=\{b_1,b_2,b_3\}$ stanowiska, a krawędź
łączy osobę ze stanowiskiem, na którym może pracować.
Załóżmy, że $N(\{a_1\})=\{b_1,b_2\}$,
$N(\{a_2\})=\{b_2\}$, $N(\{a_3\})=\{b_2,b_3\}$.
Warunek Halla jest spełniony: dla pojedynczych osób
każdy zbiór sąsiadów jest niepusty, dla par $\{a_1,a_2\},\{a_1,a_3\},\{a_2,a_3\}$
moce sąsiedztw wynoszą odpowiednio $2,3,2$,
a dla całego $A$ są trzy dostępne stanowiska.
Pusty podzbiór spełnia warunek automatycznie. Jedynym sposobem przypisania $a_2$ jest
$b_2$, po czym można wybrać $a_1b_1$ i $a_3b_3$.
Dostajemy skojarzenie doskonałe.
\end{przyklad}

\begin{center}
\begin{tikzpicture}[scale=.95,line cap=round]
 \foreach \y/\x in {1.2/1,0/2,-1.2/3}{
  \node[tgBlueVertex] (a\x) at (-2.6,\y) {$a_{\x}$};
  \node[tgRedVertex] (b\x) at (2.6,\y) {$b_{\x}$};
 }
 \draw[tgTree,line width=2.7pt] (a1)--(b1) (a2)--(b2) (a3)--(b3);
 \draw[tgGhost] (a1)--(b2) (a3)--(b2);
 \node[font=\footnotesize,text=tgNavy] at (0,-1.75)
 {Zielone: skojarzenie pokrywające wszystkie trzy osoby.};
\end{tikzpicture}
\end{center}

\begin{definicja}[System różnych reprezentantów]
Niech $I$ będzie skończonym zbiorem indeksów, a
$(A_i)_{i\in I}$ indeksowaną rodziną skończonych zbiorów.
\emph{Systemem różnych reprezentantów} tej rodziny
nazywamy iniekcję $r:I\to X$, gdzie
$X=\bigcup_{i\in I}A_i$, spełniającą $r(i)\in A_i$
dla każdego $i\in I$. Zbiory $A_i$ mogą się powtarzać:
różne indeksy wymagają różnych reprezentantów także
wtedy, gdy odpowiadają im identyczne zbiory.
\end{definicja}

\begin{wniosek}[Hall: system różnych reprezentantów]
\label{cor:tg-hall-sdr}
Niech $I$ będzie skończonym zbiorem indeksów,
a $(A_i)_{i\in I}$ rodziną skończonych zbiorów.
Istnieje system różnych reprezentantów tej rodziny
wtedy i tylko wtedy, gdy dla każdego $J\subseteq I$
zachodzi
\[
\left|\bigcup_{j\in J} A_j\right|\ge |J|.
\]
\end{wniosek}

\begin{proof}
Niech $X=\bigcup_{i\in I}A_i$. Budujemy graf dwudzielny
o częściach $L=\{0\}\times I$ i $R=\{1\}\times X$;
wprowadzamy krawędź między $(0,i)$ i $(1,x)$
dokładnie wtedy, gdy $x\in A_i$. Dzięki oznaczeniom
$0,1$ obie części są rozłączne nawet wówczas,
gdy $I\cap X\ne\varnothing$.

Jeśli $r$ jest systemem różnych reprezentantów,
krawędzie łączące $(0,i)$ z $(1,r(i))$ dla $i\in I$
tworzą skojarzenie: końce w $L$ są różne z konstrukcji,
a końce w $R$ są różne dzięki iniektywności $r$.
Skojarzenie pokrywa $L$. Odwrotnie, skojarzenie
pokrywające $L$ wyznacza dla każdego $i$ dokładnie
jednego partnera $(1,r(i))$. Istnienie krawędzi
zapewnia $r(i)\in A_i$, a rozłączność końców
krawędzi zapewnia iniektywność $r$.

Dla dowolnego $J\subseteq I$ sąsiedztwem zbioru
$\{0\}\times J$ jest dokładnie
$\{1\}\times\bigcup_{j\in J}A_j$.
Każdy podzbiór $L$ ma postać $\{0\}\times J$,
więc warunek z tezy jest równoważny warunkowi Halla
dla tego grafu. Twierdzenie~\ref{thm:tg-hall}
kończy dowód. Dla $I=\varnothing$ otrzymujemy
pustą iniekcję i warunek spełniony automatycznie.
\end{proof}

\begin{przyklad}[Powtarzające się zbiory]
Niech $I=\{1,2,3\}$ oraz $A_1=A_2=\{a,b\}$,
$A_3=\{b,c\}$. Dla pojedynczego indeksu moc
odpowiedniego zbioru wynosi $2$. Dla par indeksów
$\{1,2\},\{1,3\},\{2,3\}$ moce sum zbiorów
wynoszą odpowiednio $2,3,3$, a dla wszystkich
trzech indeksów moc sumy wynosi $3$.
Warunek Halla jest więc spełniony; przykładowym
systemem jest $r(1)=a$, $r(2)=b$, $r(3)=c$.

Gdy zastąpimy $A_3$ zbiorem $\{a,b\}$, każda
para indeksów nadal spełnia warunek Halla,
lecz dla $J=I$ suma ma tylko dwa elementy.
System trzech różnych reprezentantów nie istnieje.
\end{przyklad}

\begin{stwierdzenie}[Wzór na deficyt Halla]
\label{prop:tg-hall-def}
Niech $G=(A\sqcup B,E)$ będzie skończonym grafem
dwudzielnym. Oznaczmy przez $\nu(G)$ maksymalną
liczność skojarzenia. Wtedy
\[
\nu(G)=|A|-\max_{S\subseteq A}\bigl(|S|-|N(S)|\bigr).
\]
\end{stwierdzenie}

\begin{proof}
Niech $M$ będzie dowolnym skojarzeniem, a $S\subseteq A$.
Co najwyżej $|N(S)|$ wierzchołków $S$ może być
skojarzonych z sąsiadami, więc co najmniej
$|S|-|N(S)|$ wierzchołków $A$ pozostaje wolnych.
Zatem $|M|\le|A|-|S|+|N(S)|$; minimalizując
prawą stronę, uzyskujemy górne oszacowanie.

Niech $d:=\max_{S\subseteq A}(|S|-|N(S)|)$.
Ponieważ $S=\varnothing$ daje zero, $d\ge0$.
Dodajmy do prawej części grafu $d$ nowych
wierzchołków, każdy połączony ze wszystkimi
wierzchołkami $A$. Dla każdego niepustego $S\subseteq A$
nowe sąsiedztwo ma moc $|N(S)|+d\ge|S|$;
dla $S=\varnothing$ warunek Halla jest automatyczny.
Z twierdzenia Halla istnieje skojarzenie
pokrywające $A$ w powiększonym grafie.
Co najwyżej $d$ jego krawędzi prowadzi do nowych
wierzchołków. Po usunięciu tych krawędzi
pozostaje skojarzenie w oryginalnym grafie
o co najmniej $|A|-d$ krawędziach.
To daje nierówność odwrotną.
\end{proof}

\section{Pokrycia wierzchołkowe i twierdzenie Kőniga}
\label{sec:tg-konig}

\begin{definicja}[Pokrycie wierzchołkowe]
Dla grafu $G=(V,E)$ zbiór $C\subseteq V$ nazywamy
\emph{pokryciem wierzchołkowym}, jeśli każda krawędź
ma co najmniej jeden koniec w $C$.
Najmniejszą możliwą liczność takiego zbioru
oznaczamy przez $\tau_V(G)$.
\end{definicja}

Każde pokrycie wierzchołkowe musi zawierać przynajmniej
jeden koniec każdej krawędzi skojarzenia. Ponieważ
krawędzie skojarzenia są rozłączne w końcach,
$\nu(G)\le\tau_V(G)$ dla dowolnego grafu. Dla grafów
dwudzielnych zachodzi równość.

\begin{twierdzenie}[Kőnig]
\label{thm:tg-konig}
Niech $G=(A\sqcup B,E)$ będzie skończonym grafem
dwudzielnym. Wtedy minimalna moc pokrycia
wierzchołkowego jest równa maksymalnej liczności
skojarzenia: $\tau_V(G)=\nu(G)$.
\end{twierdzenie}

\begin{proof}
Wybieramy skojarzenie największe $M$ i oznaczamy
przez $A_0\subseteq A$ zbiór wierzchołków wolnych.
Orientujemy krawędzie spoza $M$ od $A$ do $B$,
a krawędzie $M$ od $B$ do $A$. Niech $Z_A\subseteq A$
i $Z_B\subseteq B$ będą zbiorami wierzchołków
osiągalnych z pewnego wierzchołka $A_0$ po skierowanych
ścieżkach w tym grafie. Dopuszczamy ścieżki długości
zero, więc $A_0\subseteq Z_A$. Osiągalność jest
zachowywana przy przejściu dowolnym wychodzącym
łukiem: dołączając go do ścieżki, otrzymujemy spacer,
z którego można usunąć powtórzenia wierzchołków.
Każda dodatniej długości ścieżka wychodząca z $A_0$
jest naprzemienna w wymaganym kierunku.

Żaden wierzchołek $Z_B$ nie jest wolny: w przeciwnym
razie otrzymalibyśmy ścieżkę powiększającą,
sprzecznie z twierdzeniem \ref{thm:tg-berge}.
Zatem każdy $b\in Z_B$ jest skojarzony z pewnym
$a\in Z_A$. Ponadto każdy skojarzony $a\in Z_A$
ma partnera w $Z_B$: nie należy do $A_0$, więc
każda ścieżka z $A_0$ do $a$ ma dodatnią długość,
a jej ostatni łuk z $B$ do $A$ musi należeć do $M$. Zbiory skojarzonych
wierzchołków w $Z_A$ i wszystkich wierzchołków
$Z_B$ są więc w bijekcji za pomocą krawędzi $M$.

Definiujemy $C:=(A\setminus Z_A)\cup Z_B$.
Pokażemy, że jest to pokrycie. Gdyby jakaś
krawędź $ab$ miała oba końce poza $C$,
to $a\in Z_A$ i $b\notin Z_B$.
Jeżeli $ab\notin M$, z definicji osiągalności
musielibyśmy mieć $b\in Z_B$. Jeśli zaś $ab\in M$,
to skoro $a\in Z_A$, również jego partner $b$
należy do $Z_B$. W obu przypadkach sprzeczność.

Każda krawędź $M$ ma dokładnie jeden koniec w $C$:
partnerzy z $Z_A$ i $Z_B$ przechodzą do $C$ tylko
po stronie $B$, a partnerzy poza nimi tylko po
stronie $A$. Żaden wierzchołek wolny nie należy
do $C$ (wolne w $A$ są w $Z_A$, a wolnych
w $Z_B$ nie ma). W konsekwencji $|C|=|M|$.
Ponieważ zawsze $\nu(G)\le\tau_V(G)$, a znaleźliśmy
pokrycie o mocy $|M|=\nu(G)$, otrzymujemy równość.
\end{proof}

\begin{uwaga}
Twierdzenie Bergego charakteryzuje największe skojarzenia
w dowolnym skończonym grafie prostym przez brak ścieżek
powiększających. W grafach dwudzielnych twierdzenie Halla
i wzór na deficyt opisują przeszkodę przez sąsiedztwa
zbiorów, a twierdzenie Kőniga wyraża liczność największego
skojarzenia za pomocą najmniejszego pokrycia wierzchołkowego. W dalszej części skryptu
podobne idee będą wyrażane przez rząd macierzy,
przekroje i przepływy, a później przez grupy homologii.
\end{uwaga}

\section{Ścieżki i cykle Hamiltona}
\label{sec:tg-hamilton}

\begin{definicja}[Hamiltonowskość]
Niech $G$ będzie skończonym grafem prostym.
\emph{Ścieżką Hamiltona} nazywamy ścieżkę zawierającą
wszystkie wierzchołki $G$. \emph{Cyklem Hamiltona}
nazywamy cykl zawierający wszystkie wierzchołki $G$.
Graf z cyklem Hamiltona nazywamy \emph{hamiltonowskim}.
Zgodnie z przyjętą definicją cyklu prostego taki graf
ma co najmniej trzy wierzchołki. Graf $K_1$ ma ścieżkę
Hamiltona długości zero, lecz graf zerowy nie ma
żadnej ścieżki.
\end{definicja}

Obecność cyklu Eulera i Hamiltona to dwa różne warunki.
Cykl Eulera przechodzi przez wszystkie krawędzie,
a cykl Hamiltona przez wszystkie wierzchołki.
Cykl $C_n$ spełnia oba warunki. Dwa trójkąty o jednym
wspólnym wierzchołku mają cykl Eulera, ale nie mają
cyklu Hamiltona: usunięcie wspólnego wierzchołka
rozspaja graf, podczas gdy usunięcie dowolnego
wierzchołka z cyklu Hamiltona pozostawia ścieżkę
zawierającą wszystkie pozostałe wierzchołki.

\begin{center}
\begin{tikzpicture}[scale=.95,line cap=round]
 \foreach \x/\y/\n in {-2/0/a,-1/1.25/b,0/0/c,1/1.25/d,2/0/e}
   \node[tgBlueVertex] (\n) at (\x,\y) {$\n$};
 \draw[tgTree] (a)--(b)--(c)--(a);
 \draw[tgAccent] (c)--(d)--(e)--(c);
 \node[font=\footnotesize,text=tgNavy] at (0,-.65)
   {Euler: tak. Hamilton: nie (wierzchołek rozcinający $c$).};
\end{tikzpicture}
\end{center}

\begin{stwierdzenie}[Konieczny warunek hamiltonowskości]
\label{prop:tg-hamilton-components}
Niech $G=(V,E)$ będzie skończonym grafem prostym
z cyklem Hamiltona. Dla każdego niepustego
$S\subseteq V$ zachodzi $c(G-S)\le |S|$, gdzie
$c(H)$ oznacza liczbę składowych spójnych grafu $H$
(i $c(H)=0$, gdy $H$ jest grafem zerowym).
\end{stwierdzenie}

\begin{proof}
Niech $C$ będzie cyklem Hamiltona w $G$.
Jeśli $S=V$, to $G-S$ jest grafem zerowym
i nierówność jest oczywista. Niech więc
$\varnothing\ne S\subsetneq V$.
Ustalmy kierunek obiegu $C$ i rozpocznijmy go
w dowolnym wierzchołku z $S$. Po usunięciu $S$
pozostają niepuste bloki kolejnych wierzchołków
cyklu spoza $S$. Każdy taki blok jest ścieżką
(być może długości zero) i bezpośrednio poprzedza
go w obiegu wierzchołek z $S$. Dwa różne bloki
nie mogą mieć tego samego poprzednika, ponieważ
każdy wierzchołek występuje w obiegu dokładnie raz.
Składowe $C-S$ są właśnie tymi blokami, więc
$c(C-S)\le |S|$.

Graf $C-S$ jest podgrafem rozpinającym $G-S$:
oba mają ten sam zbiór wierzchołków $V\setminus S$,
a każda krawędź pierwszego należy do drugiego.
Każda składowa $C-S$ zawiera się zatem w pewnej
składowej $G-S$, a każda składowa $G-S$ zawiera
co najmniej jedną składową $C-S$.
Stąd $c(G-S)\le c(C-S)\le |S|$.
\end{proof}

\begin{przyklad}[Brak wierzchołków rozcinających nie wystarcza]
Rozważmy graf pełny dwudzielny $G=K_{2,3}$,
o częściach $A=\{a_1,a_2\}$ i $B=\{b_1,b_2,b_3\}$.
Usunięcie dowolnego wierzchołka daje $K_{1,3}$
lub $K_{2,2}$, a oba te grafy są spójne.
Istotnie, w każdym z nich wierzchołki z różnych
części są sąsiadami, a dwa wierzchołki z tej samej
części łączy ścieżka długości dwa przez dowolny
wierzchołek drugiej części. Graf $G$ nie ma więc
wierzchołków rozcinających.
Jednak dla $S=A$ graf $G-S$ składa się z trzech
izolowanych wierzchołków, czyli
$c(G-S)=3>2=|S|$.
Na mocy stwierdzenia~\ref{prop:tg-hamilton-components}
graf $G$ nie ma cyklu Hamiltona.
\end{przyklad}

\begin{twierdzenie}[Dirac]
\label{thm:tg-dirac}
Niech $G$ będzie skończonym grafem prostym o
$n\ge3$ wierzchołkach i minimalnym stopniu
$\delta(G)\ge n/2$. Wtedy $G$ ma cykl Hamiltona.
\end{twierdzenie}

\begin{proof}
Najpierw zauważmy, że $G$ jest spójny.
Gdyby miał co najmniej dwie składowe, w każdej
z nich liczba wierzchołków byłaby co najmniej
$\delta(G)+1\ge n/2+1$. Suma liczebności dwóch
składowych przekroczyłaby $n$, sprzeczność.

Wybierzmy ścieżkę o największej liczbie wierzchołków:
$P=v_1,v_2,\ldots,v_k$. Ponieważ każdy wierzchołek
ma co najmniej dwa różne sąsiady ($n\ge3$),
istnieje ścieżka o trzech wierzchołkach, więc $k\ge3$. Wszystkich sąsiadów obu
końców $v_1,v_k$ musi zawierać $P$, bo w przeciwnym
razie przedłużylibyśmy ścieżkę.
Dla $i=1,\ldots,k-1$ rozważmy dwa warunki:
$v_1v_{i+1}\in E$ oraz $v_iv_k\in E$.
Pierwszy jest spełniony dla dokładnie $d(v_1)$ indeksów,
a drugi dla dokładnie $d(v_k)$ indeksów.
Ponieważ $d(v_1)+d(v_k)\ge n>k-1$, obydwa
zbiory indeksów muszą mieć punkt wspólny.
Wybierzmy takie $i$.

Wtedy ciąg
$v_1,v_2,\ldots,v_i,v_k,v_{k-1},\ldots,v_{i+1},v_1$
stanowi cykl prosty: krawędzie wewnątrz fragmentów
pochodzą z $P$, przejścia $v_i v_k$ i $v_{i+1}v_1$
wynikają z wyboru $i$, a rozłączne listy
$v_1,\ldots,v_i$ i $v_k,\ldots,v_{i+1}$ zawierają
łącznie każdy wierzchołek $P$ dokładnie raz.
Ponieważ $k\ge3$, jest to cykl w przyjętym sensie.

Gdyby $k<n$, spójność $G$ zapewniałaby istnienie
krawędzi między tym cyklem a pewnym wierzchołkiem
$w\notin V(P)$ (bierzemy pierwszą krawędź ścieżki
wychodzącej poza cykl). Niech $x$ będzie wierzchołkiem cyklu sąsiadującym z $w$.
Usuwamy jedną z dwóch krawędzi cyklu incydentnych z $x$.
Pozostaje ścieżka o $k$ wierzchołkach, której jednym
końcem jest $x$. Przechodząc nią od drugiego końca do
$x$, a następnie krawędzią $xw$, otrzymujemy ścieżkę
o $k+1$ różnych wierzchołkach.
Byłaby dłuższa od $P$, sprzeczność.
Zatem $k=n$, a znaleziony cykl jest Hamiltona.
\end{proof}

\begin{twierdzenie}[Ore]
\label{thm:tg-ore}
Niech $G$ będzie skończonym grafem prostym o
$n\ge3$ wierzchołkach. Jeżeli dla każdej pary
różnych niesąsiadujących wierzchołków $u,v$
zachodzi $d(u)+d(v)\ge n$, to $G$ jest hamiltonowski.
\end{twierdzenie}

\begin{proof}
Wykażemy najpierw pomocniczą zasadę: jeśli $u,v$
są niesąsiadujące i $d(u)+d(v)\ge n$, to
dodanie krawędzi $uv$ nie zmienia odpowiedzi
na pytanie o istnienie cyklu Hamiltona.
Jeden kierunek jest oczywisty: cykl Hamiltona
w $G$ pozostaje cyklem w $G+uv$.

Załóżmy, że w $G+uv$ istnieje cykl Hamiltona
używający nowej krawędzi $uv$ (inaczej już jest
cyklem w $G$). Po usunięciu tej krawędzi otrzymujemy
ścieżkę Hamiltona $P=x_1,x_2,\ldots,x_n$ w $G$
z $x_1=u$ i $x_n=v$. Dla $i=1,\ldots,n-1$
rozważmy warunki $ux_{i+1}\in E$ i $vx_i\in E$.
Liczby indeksów spełniających te warunki to
odpowiednio $d(u)$ i $d(v)$ (ponieważ $u,v$
nie sąsiadują). Ponieważ ich suma wynosi co
najmniej $n>n-1$, istnieje wspólny indeks $i$.
Wtedy
$u=x_1,x_2,\ldots,x_i,v=x_n,x_{n-1},\ldots,x_{i+1},u$
jest cyklem Hamiltona w $G$: oba przejścia między
fragmentami ścieżki wynikają z wyboru wspólnego indeksu,
pozostałe krawędzie należą do $P$, a oba fragmenty
zawierają łącznie wszystkie $n$ różnych wierzchołków.

Dodawajmy teraz kolejno wszystkie brakujące
krawędzie grafu $G$. Przy każdym kroku suma stopni
końców dodawanej krawędzi jest co najmniej $n$,
bo stopnie mogą tylko wzrastać. Z udowodnionej
zasady dodawanie żadnej z krawędzi nie zmienia
hamiltonowskości. Ostatecznie otrzymujemy $K_n$,
który ma cykl Hamiltona. Zatem również $G$ go ma.
\end{proof}

\begin{przyklad}[Porównanie kryteriów Diraca i Ore'a]
Graf $K_{3,3}$ ma $n=6$ i wszystkie stopnie równe trzy,
więc spełnia warunek Diraca i jest hamiltonowski.
Z kolei cykl $C_5$ ma stopnie równe dwa; nie spełnia
warunku Diraca, a dla każdej pary niesąsiadujących
wierzchołków suma stopni wynosi $4<5$, więc nie spełnia
również warunku Ore'a. Sam jest jednak cyklem Hamiltona.
Warunki Diraca i Ore'a są więc wystarczające, niekonieczne.

Istnieją także grafy, dla których działa kryterium Ore'a,
ale kryterium Diraca już nie. Rozważmy graf $K_5$
po usunięciu dwóch krawędzi o wspólnym końcu $a$.
Jego stopnie to $2,3,3,4,4$, zatem
$\delta(G)=2<5/2$. Jedyne pary niesąsiadujące
to końce dwóch usuniętych krawędzi; dla każdej
suma stopni wynosi $2+3=5=n$.
Z twierdzenia Ore'a wynika hamiltonowskość.
\end{przyklad}

\begin{center}
\begin{tikzpicture}[scale=.92,line cap=round]
  \foreach \ang/\name in {90/a,18/b,-54/c,-126/d,162/e}{
    \node[tgBlueVertex] (\name) at (\ang:1.7) {$\name$};
  }
  \foreach \u/\v in {a/d,a/e,b/c,b/d,b/e,c/d,c/e,d/e}
    \draw[tgEdge] (\u)--(\v);
  \node[font=\footnotesize,text=tgNavy] at (0,-2.2)
     {$K_5$ bez dwóch krawędzi przy $a$: wariant z warunkiem Ore'a.};
\end{tikzpicture}
\end{center}

\section*{Od kombinatoryki do algebry i topologii}
\addcontentsline{toc}{section}{Od kombinatoryki do algebry i topologii}

Dotychczas rozpoznawaliśmy cykle, drzewa, rozłączne ścieżki
i skojarzenia metodami kombinatorycznymi. W najbliższym bloku
(od rozdziału \ref{ch:tg-bfs-dfs} do rozdziału \ref{ch:tg-expanders})
rozwiniemy algorytmy, probabilistyczną teorię grafów i własności
spektralne. Następnie, w rozdziale \ref{ch:tg-algebra},
wprowadzimy macierze incydencji i przestrzenie łańcuchów.
Pozwoli to zrozumieć cykle grafów również jako elementy grup homologii.


% ============================================================
% ITERACJA 3. Algebra, homologia, grafy jako CW, laplasjan.
% Tekst wymaga preambuły istniejącego skryptu.
% ============================================================

